\documentclass{article}
\usepackage{amsmath, amssymb, amsthm}
\usepackage{hyperref}
\usepackage{geometry}
\geometry{margin=1in}

\title{Erd\H{o}s Problem \#830}
\date{Last updated: 2025-08-31}
\author{Source: erdosproblems.com}

\begin{document}
\maketitle

\noindent\textbf{Status:} OPEN \quad \textbf{No prize}

\noindent\textbf{Tags:} number theory

\medskip

\noindent\textbf{Problem Statement:}

We say that $a,b\in \mathbb{N}$ are an amicable pair if $\sigma(a)=\sigma(b)=a+b$. Are there infinitely many amicable pairs? If $A(x)$ counts the number of amicable $1\leq a\leq b\leq x$ then is it true that\[A(x)&gt;x^{1-o(1)}?\]




For example $220$ and $284$. Erd\H{o}s \cite{Er55b} proved that $A(x)=o(x)$, and Pomerance \cite{Po81} improved this to\[A(x) \leq x \exp(-(\log x)^{1/3})\]and later \cite{Po15} to\[A(x) \leq x \exp(-(\tfrac{1}{2}+o(1))(\log x\log\log x)^{1/2}).\]This is problem B4 in Guy's collection \cite{Gu04}.




References

[Er55b] Erd\H{o}s, P., On amicable numbers . Publ. Math. Debrecen (1955), 108-111.

[Gu04] Guy, Richard K., Unsolved problems in number theory . (2004), xviii+437.

[Po15] Pomerance, Carl, On amicable numbers . (2015), 321-327.

[Po81] Pomerance, Carl, On the distribution of amicable numbers. {II} . J. Reine Angew. Math. (1981), 183-188.

\noindent\textbf{Related OEIS sequences:} A259180


\bigskip
\noindent\small{Source: \url{https://www.erdosproblems.com/830}}
\end{document}
